% ═══════════════════════════════════════════════════════════
% §4 Methods
% ═══════════════════════════════════════════════════════════
\section{Methods}
\label{sec:methods}

\subsection{Gray-code encoding of biological alphabets}

A reflected binary Gray code assigns to each non-negative integer $n$ the
value $g(n) = n \oplus (n \gg 1)$, where $\oplus$ denotes bitwise exclusive-or.
The defining property is that consecutive integers map to codewords differing
in exactly one bit position~\cite{gray1953}.  For nucleotide sequences we adopt
the standard two-bit encoding A=0, C=1, G=3, T/U=2.  For amino acids we assign
each residue its He~2012 integer code~\cite{he2012} under an ordering that
groups residues into seven physicochemical classes (aliphatic: A,I,L,V;
aromatic: F,Y,W; sulfur: M,C; structure: P,G; polar: S,T,N,Q; negative: D,E;
positive: H,K,R).  The five-bit Gray image $g(\mathrm{code}(a))$ places each
residue on the five-dimensional hypercube $Q_5$ such that within-group pairs
tend toward small Hamming distance.  Among the $\binom{20}{2} = 190$ unordered
residue pairs, exactly 40 lie at Gray distance~1 (verified exhaustively in
Lean~4).  Alignment gap characters are mapped to a twenty-first state
(code~20) that is excluded from all statistical calculations but retained in
position arrays to preserve column identity.

\subsection{Karnaugh-map construction and Quine--McCluskey minimisation}

For a pair of alignment columns $(i,j)$, a ternary Karnaugh map assigns each
cell $(a,b)$ of a $32\times32$ grid one of three values based on the joint
frequency of consensus residues across the alignment: ON if $(a,b)$ is an
observed mutation, OFF if never observed, or don't-care if it is the reference
(majority) pair.  Rows and columns with indices 20--31 correspond to unused
Gray codewords and receive don't-care status.  The map is flattened row-major
into a 1024-element truth table for minimisation input.

The Quine--McCluskey algorithm~\cite{quine1952,mccluskey1956} finds all prime
implicants by iterative pairwise merging of minterms differing in one bit,
with don't-care cells included in the initial minterm set alongside on-set
cells.  Essential prime implicants cover on-set cells matched by no other
implicant; a greedy cover completes the solution.  Every trained circuit is
asserted sound (no implicant matches any OFF cell) and complete (every ON cell
is covered) before evaluation, properties formalised as theorems in Lean~4.

\subsection{Information-theoretic measures}

Shannon entropy per column is computed after excluding gap characters.
Mutual information between columns $(i,j)$ uses the standard formula with
20-state joint distribution over canonical residues.  Perplexity ratio is
defined as the marginal perplexity of column $j$ divided by the mean
conditional perplexity given residue $a$ at column $i$, averaged over residues
with $\geq 5$ observations.  Coupling constants use the log-odds ratio
$J_{ij}(a,b) = \ln(P_{ij}(a,b)/(P_i(a)P_j(b)))$.

\subsection{Direct coupling analysis}

Mean-field DCA follows Morcos et al.\cite{morcos2011}: sequence reweighting at
$\theta = 0.2$, pseudocount $\lambda = 0.5$, covariance inversion over
$(q{-}1)L \times (q{-}1)L$ matrix with gap state removed, Frobenius-norm
scoring, and APC correction~\cite{dunn2008}.  Alignments deeper than 12{,}000
sequences are uniformly subsampled.  Regularisation uses a ridge of
$10^{-2}$ added to the diagonal before inversion.

\subsection{Dataset and evaluation}

The PSICOV150 benchmark~\cite{jones2012} provides 150 protein domains with
deep MSAs and repaired native structures.  The number of alignment columns
equals the number of C-$\alpha$ atoms in each PDB file, verified for 149/150
targets.  Native contacts use C$\beta$--C$\beta < 8$~\AA{} (C$\alpha$ for
glycine) with separation $\geq 6$.  Published PSICOV contact predictions from
the supplementary data serve as literature baseline.

Cross-protein evaluation uses protein-level cross-validation with folds
stratified by length quintile (seed~42).  All model parameters derive
exclusively from training proteins.  Ranked metrics include precision@L/5,
precision@L/2, precision@L, and AUC via Mann--Whitney $U$ with mid-rank tie
handling.  Enrichment factors express precision relative to per-dataset base
contact rate.

\subsection{Formal verification in Lean~4}

All load-bearing definitions are formalised in Lean~4~\cite{lean4} across two
packages totalling 118 theorems.  Key results include Gray-code consecutive
adjacency, Hamming-distance symmetry, He-code pair counts at each distance,
separation-bin monotonicity, level-1 cell-index injectivity, implicant-cover
completeness and soundness, fully-fixed implicant uniqueness on the padded
grid, and padding safety (border codes cannot leak into decoded rules).
Every theorem depends only on \texttt{propext}, \texttt{Quot.sound},
\texttt{Classical.choice}, and the \texttt{native\_decide} evaluation axiom;
none depends on \texttt{sorryAx}.
